\begin{abstract}
% This paper presents a lightweight browser-based 3D engine designed to support the development and execution of interactive educational scenarios without requiring native runtime installation. The engine is needed because existing approaches to delivering 3D learning content through the browser force developers to choose between heavyweight exported runtimes with long loading times and low-level graphics libraries that demand specialized rendering expertise. The proposed solution provides a high-level component model over a hidden rendering backend, enabling educational platform developers to build interactive 3D scenarios with reduced engineering effort while maintaining runtime efficiency on consumer hardware. The paper describes the key architectural choices behind the engine and introduces three targeted optimization strategies addressing asset memory consumption, rendering initialization latency, and per-frame computational overhead. The evaluation covered 87 configurations across three benchmark scenes, using three graphics adapters on two machines. Each strategy is reported against the metric it is capable of moving. The results indicate that the approach is practical for browser-delivered educational 3D content, and that which hardware benefits from which optimization does not follow a simple low-end to high-end ordering.
This paper presents a lightweight browser-based 3D engine designed to support the development and execution of interactive educational scenarios without requiring native runtime installation. Existing approaches to delivering 3D learning content through the browser force developers to choose between heavyweight exported runtimes with long loading times and low-level graphics libraries that demand specialized rendering expertise. The proposed solution provides a high-level component model over a hidden rendering backend, enabling educational platform developers to build interactive 3D scenarios with reduced engineering effort while maintaining runtime efficiency on consumer hardware. The paper describes the key architectural choices behind the engine and introduces three targeted optimization strategies addressing asset memory consumption, rendering initialization latency, and per-frame computational overhead. The evaluation covered 87 configurations across three benchmark scenes, using three graphics adapters on two machines. GPU-compressed textures reduce the cost of the first rendered frame by 56.7\% to 69.5\%; shader pre-compilation is effective on the discrete adapters but not on the integrated one; and dirty-flag transform propagation trades a per-phase speed-up for a frame-level cost in the workload measured. The results indicate that the approach is practical for browser-delivered educational 3D content, and that the benefit of a given optimization does not follow a simple low-end to high-end ordering.

\keywords{WebGL \and game engine architecture \and interactive learning \and 3D visualization \and performance optimization \and technology-enhanced learning}

\end{abstract}